\documentclass[11pt,a4paper]{article}
\usepackage[utf8]{inputenc}
\usepackage{amsmath,amssymb}
\usepackage{geometry}
\geometry{margin=2.5cm}
\usepackage{listings}
\usepackage{xcolor}
\usepackage{hyperref}

\lstset{
  language=C++,
  basicstyle=\ttfamily\small,
  keywordstyle=\color{blue}\bfseries,
  commentstyle=\color{green!60!black},
  stringstyle=\color{red!70!black},
  numbers=left,
  numberstyle=\tiny\color{gray},
  numbersep=8pt,
  frame=single,
  breaklines=true,
  tabsize=4,
  showstringspaces=false
}

\title{IOI 1998 -- Polygon}
\author{Solution}
\date{}

\begin{document}
\maketitle

\section{Problem Statement}

A polygon-shaped game has $n$ vertices labeled with integers and $n$ edges labeled
with \texttt{+} (addition) or \texttt{*} (multiplication). The vertices and edges
alternate around the polygon: edge $i$ connects vertex $i$ to vertex $i+1$
(with edge $n$ connecting vertex $n$ to vertex 1).

The game proceeds in two phases:
\begin{enumerate}
  \item Remove one edge, converting the polygon into a chain of $n$ vertices and $n-1$ edges.
  \item Repeatedly merge two adjacent vertices via their connecting edge: replace
        the pair with a single vertex whose value is the result of the edge's operation
        applied to the two vertex values. Continue until one vertex remains.
\end{enumerate}

Find the maximum possible value of the final vertex, considering all choices of
which edge to remove and all merge orders.

\medskip
\noindent\textbf{Constraints:} $3 \le n \le 50$, vertex values in $[-32768, 32767]$.

\section{Solution Approach}

\subsection{Circular Interval DP}

The problem is analogous to \textbf{optimal matrix chain multiplication} with the
added complication that multiplication of negative numbers can turn a minimum into
a maximum. We track both the maximum and minimum achievable values for each sub-interval.

\subsection{Handling the Circle}

Double the sequence: vertices $v_0, \ldots, v_{n-1}, v_0, \ldots, v_{n-1}$ (and
similarly for edges). Then every contiguous sub-chain of length $n$ corresponds to
removing one particular edge.

\subsection{DP with Min and Max}

For interval $[i, j]$ of the doubled sequence, define:
\begin{align*}
  \mathrm{maxDP}[i][j] &= \text{maximum value achievable from vertices } i \text{ through } j, \\
  \mathrm{minDP}[i][j] &= \text{minimum value achievable from vertices } i \text{ through } j.
\end{align*}

\textbf{Base case:} $\mathrm{maxDP}[i][i] = \mathrm{minDP}[i][i] = v_i$.

\textbf{Recurrence:} For each split point $k \in [i, j{-}1]$, let the edge between
positions $k$ and $k+1$ have operation $\mathrm{op}$:
\begin{itemize}
  \item If $\mathrm{op} = $ \texttt{+}:
    \begin{align*}
      \mathrm{maxDP}[i][j] &\gets \max\bigl(\mathrm{maxDP}[i][k] + \mathrm{maxDP}[k{+}1][j]\bigr), \\
      \mathrm{minDP}[i][j] &\gets \min\bigl(\mathrm{minDP}[i][k] + \mathrm{minDP}[k{+}1][j]\bigr).
    \end{align*}
  \item If $\mathrm{op} = $ \texttt{*}: consider all four products of min/max from each half,
    since multiplying two negatives yields a positive:
    \begin{align*}
      \text{candidates} = \{&\mathrm{maxDP}[i][k] \cdot \mathrm{maxDP}[k{+}1][j],\;
                             \mathrm{maxDP}[i][k] \cdot \mathrm{minDP}[k{+}1][j], \\
                            &\mathrm{minDP}[i][k] \cdot \mathrm{maxDP}[k{+}1][j],\;
                             \mathrm{minDP}[i][k] \cdot \mathrm{minDP}[k{+}1][j]\}.
    \end{align*}
\end{itemize}

\subsection{Edge Indexing}

In the doubled array, the edge between positions $k$ and $k+1$ is edge
$(k + 1) \bmod n$ from the original polygon. This is because the input format gives
edge $i$ \emph{before} vertex $i$: the $i$-th input line contains edge $i$'s operation
followed by vertex $i$'s value, where edge $i$ connects vertex $i{-}1$ to vertex $i$
(cyclically). So the edge between positions $k$ and $k+1$ in our 0-indexed array
corresponds to edge $(k+1) \bmod n$.

\subsection{Answer}

The answer is $\max_{0 \le i < n} \mathrm{maxDP}[i][i + n - 1]$.

\section{C++ Solution}

\begin{lstlisting}
#include <cstdio>
#include <cstring>
#include <algorithm>
#include <climits>
using namespace std;

typedef long long ll;
const ll NEG_INF = -1e18;
const ll POS_INF = 1e18;

int n;
ll val[105];     // vertex values (doubled)
char op[105];    // edge operations (doubled): 't' for +, 'x' for *

ll maxDP[105][105], minDP[105][105];

int main() {
    scanf("%d", &n);

    for (int i = 0; i < n; i++) {
        char c;
        int v;
        scanf(" %c %d", &c, &v);
        op[i] = c;
        val[i] = v;
    }

    // Double the arrays for circular handling
    for (int i = 0; i < n; i++) {
        val[i + n] = val[i];
        op[i + n] = op[i];
    }

    int N = 2 * n;

    // Initialize DP
    for (int i = 0; i < N; i++)
        for (int j = 0; j < N; j++) {
            maxDP[i][j] = NEG_INF;
            minDP[i][j] = POS_INF;
        }

    for (int i = 0; i < N; i++) {
        maxDP[i][i] = val[i];
        minDP[i][i] = val[i];
    }

    // Fill DP for increasing interval lengths
    for (int len = 2; len <= n; len++) {
        for (int i = 0; i + len - 1 < N; i++) {
            int j = i + len - 1;
            for (int k = i; k < j; k++) {
                // Edge between position k and k+1 is op[(k+1) % n]
                char o = op[(k + 1) % n];

                if (o == 't') {
                    maxDP[i][j] = max(maxDP[i][j],
                        maxDP[i][k] + maxDP[k+1][j]);
                    minDP[i][j] = min(minDP[i][j],
                        minDP[i][k] + minDP[k+1][j]);
                } else { // multiplication
                    ll cands[4] = {
                        maxDP[i][k] * maxDP[k+1][j],
                        maxDP[i][k] * minDP[k+1][j],
                        minDP[i][k] * maxDP[k+1][j],
                        minDP[i][k] * minDP[k+1][j]
                    };
                    for (int c = 0; c < 4; c++) {
                        maxDP[i][j] = max(maxDP[i][j], cands[c]);
                        minDP[i][j] = min(minDP[i][j], cands[c]);
                    }
                }
            }
        }
    }

    // Find maximum over all starting positions (each = removing one edge)
    ll ans = NEG_INF;
    for (int i = 0; i < n; i++)
        ans = max(ans, maxDP[i][i + n - 1]);

    printf("%lld\n", ans);

    // Print which edge removals achieve the maximum (1-indexed)
    for (int i = 0; i < n; i++)
        if (maxDP[i][i + n - 1] == ans)
            printf("%d ", i + 1);
    printf("\n");

    return 0;
}
\end{lstlisting}

\section{Correctness}

\begin{itemize}
  \item We track both $\mathrm{maxDP}$ and $\mathrm{minDP}$ because multiplication with negative
        numbers can invert the ordering. All four products of extrema are considered.
  \item The doubled-array technique correctly enumerates all $n$ ways to break the circle.
  \item Using \texttt{long long} prevents overflow: vertex values up to $32767$ and
        up to 50 multiplications could produce values up to $32767^{50}$, but in practice
        the intermediate values stay within \texttt{long long} range due to the alternation
        of additions and multiplications, and the constraint structure.
\end{itemize}

\section{Complexity Analysis}

\begin{itemize}
  \item \textbf{Time complexity:} $O(n^3)$. There are $O(n^2)$ intervals in the doubled
        array (only those of length $\le n$), each with $O(n)$ split points.
        For $n = 50$: about $125{,}000$ operations.
  \item \textbf{Space complexity:} $O(n^2)$ for the DP tables.
\end{itemize}

\end{document}
